\documentclass[a4paper,11pt]{article}

\usepackage{amsmath}
\usepackage{amssymb}
\usepackage{booktabs}
\usepackage{hyperref}
\usepackage{geometry}
\usepackage{titlesec}
\geometry{top=0.7in, right=0.75in, left=0.75in, bottom=0.7in}

\titlespacing*{\section}{0pt}{8pt}{6pt}
\titlespacing*{\subsection}{0pt}{6pt}{2pt}
\linespread{1}

\title{Detailed Appendix\\[0.2em]\large Full derivations for Q4 (cap) and Q5 (RTM bargain)}
\author{Alessandro Caggia (3178115), Stefano De Bianchi (3154360), Eduard Hartel (3378615), Francesco Mastidoro (3323853)}
\date{May 2026}

\begin{document}
\maketitle

\noindent\emph{Companion document to \texttt{problem1\_v2.pdf}.} The main paper presents streamlined versions of \S A.1 (primitives, FOCs, bliss points) and \S A.2 (firm's constrained programme); the unabridged derivations, with full KKT machinery and case-by-case argument, are reproduced below. Sections \S A.3--\S A.5 and Appendix B (\S B.1--\S B.7) are reproduced verbatim from the main paper.

% =========================================================================
\section{Formal derivation of the Q4 cap results}
% =========================================================================

\subsection{Primitives, FOCs, and the three bliss points}
The worker and firm payoffs are
\[
V_w(h)=(1-\tau)wh-\tfrac{\alpha}{2}h^2,
\qquad
V_f(h)=Ah-\tfrac{\beta}{2}h^2-wh.
\]
\textbf{Worker FOC.} $V_w'(h)=(1-\tau)w-\alpha h=0\;\Longrightarrow\;h^{*}=(1-\tau)w/\alpha$ (the worker's preferred hours absent any constraint).

\textbf{Firm FOC.} $V_f'(h)=A-\beta h-w=0\;\Longrightarrow\;h^{\mathrm{firm}}=(A-w)/\beta$ (the firm's unconstrained labour demand at wage $w$).

\textbf{Joint-surplus FOC.} Adding the two payoffs the wage transfer cancels and
\[
W(h)=V_w(h)+V_f(h)=Ah-\tau wh-\tfrac{\alpha+\beta}{2}h^{2},\qquad W'(h)=A-\tau w-(\alpha+\beta)h=0\;\Longrightarrow\;h^{\mathrm{soc}}=\frac{A-\tau w}{\alpha+\beta}.
\]
At the calibration $(A,w,\tau,\alpha,\beta)=(2,1,0.4,2,2)$: $h^{*}=0.30$, $h^{\mathrm{soc}}=0.40$, $h^{\mathrm{firm}}=0.50$, so the firm-power ordering $h^{*}<h^{\mathrm{soc}}<h^{\mathrm{firm}}$ holds.

\subsection{Firm's constrained programme under a cap}
In the firm-power baseline the firm chooses $h$ to maximise its own surplus subject to (i) the regulatory cap $h\le\bar h$ and (ii) the worker's individual-rationality (IR) constraint $V_w(h)\ge V_w^{0}$, with the wage chosen so that IR binds.

\textbf{Why IR enters the programme.} Without IR the firm-power problem is degenerate (the firm could set $w$ arbitrarily low and extract unbounded rents); IR pins the wage at the lowest level that satisfies the worker's outside option $V_w^{0}$, closing the model.

The Lagrangian is
\[
\mathcal{L}(h,\mu,\lambda)=V_f(h)-\mu(h-\bar h)+\lambda\bigl(V_w(h)-V_w^{0}\bigr),\qquad \mu,\lambda\ge 0.
\]
The FOC and KKT slackness conditions read
\[
V_f'(h)+\lambda V_w'(h)=\mu,\qquad \mu(h-\bar h)=0,\qquad \lambda(V_w(h)-V_w^{0})=0.
\]
The wage is set so that IR binds, which fixes $\lambda$; the inner choice of $h$ then collapses to $\max_h V_f(h)$ s.t.\ $h\le\bar h$. Two cases:
\begin{itemize}\itemsep0pt
\item \textbf{Cap slack} ($\bar h\ge h^{\mathrm{firm}}$): $\mu=0$, $h^{\mathrm{eq}}=h^{\mathrm{firm}}$.
\item \textbf{Cap binding} ($\bar h<h^{\mathrm{firm}}$): $h^{\mathrm{eq}}=\bar h$ and $\mu=V_f'(\bar h)=A-w-\beta\bar h>0$ (the shadow price of the cap is the firm's marginal value of an extra hour at the constrained equilibrium).
\end{itemize}
Combined: $h^{\mathrm{eq}}(\bar h)=\min\{\bar h,\,h^{\mathrm{firm}}\}$. All subsequent surpluses are evaluated at this $h^{\mathrm{eq}}$ relative to the no-cap baseline $h^{\mathrm{firm}}$.

\subsection{Quadratic identities for the three surpluses}
Any concave quadratic $g(h)$ with bliss point $h^{\dagger}$ and curvature $-\kappa$ satisfies
\begin{equation}
g(h)-g(h^{\dagger})=-\tfrac{\kappa}{2}(h-h^{\dagger})^{2}. \label{eq:quad-id}
\end{equation}
Apply \eqref{eq:quad-id} to $V_f$ (bliss $h^{\mathrm{firm}}$, curvature $\beta$) and to $W$ (bliss $h^{\mathrm{soc}}$, curvature $\alpha+\beta$):
\begin{align}
\Delta V_f(\bar h) &\equiv V_f(\bar h)-V_f(h^{\mathrm{firm}}) = -\tfrac{\beta}{2}\bigl(h^{\mathrm{firm}}-\bar h\bigr)^{2}\le 0, \label{eq:dVf}\\
\Delta W(\bar h) &\equiv W(\bar h)-W(h^{\mathrm{firm}})
   = \bigl[W(\bar h)-W(h^{\mathrm{soc}})\bigr] - \bigl[W(h^{\mathrm{firm}})-W(h^{\mathrm{soc}})\bigr] \nonumber\\
   &= -\tfrac{\alpha+\beta}{2}\Bigl[(h^{\mathrm{soc}}-\bar h)^{2}-(h^{\mathrm{soc}}-h^{\mathrm{firm}})^{2}\Bigr]. \label{eq:dW}
\end{align}
For $\Delta V_w$ the reference point $h^{\mathrm{firm}}$ is \emph{not} the worker's bliss, so we expand directly:
\begin{align}
\Delta V_w(\bar h)
&= V_w(\bar h)-V_w(h^{\mathrm{firm}}) \nonumber\\
&= (1-\tau)w(\bar h-h^{\mathrm{firm}})-\tfrac{\alpha}{2}\bigl(\bar h^{2}-(h^{\mathrm{firm}})^{2}\bigr) \nonumber\\
&= (\bar h-h^{\mathrm{firm}})\Bigl[(1-\tau)w-\tfrac{\alpha}{2}(\bar h+h^{\mathrm{firm}})\Bigr], \label{eq:dVw}
\end{align}
using the factorisation $\bar h^{2}-(h^{\mathrm{firm}})^{2}=(\bar h-h^{\mathrm{firm}})(\bar h+h^{\mathrm{firm}})$.

\subsection{Welfare-improving band: derivation of the optimality condition}
For $\bar h<h^{\mathrm{firm}}$, \eqref{eq:dW} gives
\[
\Delta W(\bar h)>0 \;\Longleftrightarrow\; (h^{\mathrm{soc}}-\bar h)^{2}<(h^{\mathrm{soc}}-h^{\mathrm{firm}})^{2}
\;\Longleftrightarrow\; |h^{\mathrm{soc}}-\bar h|<h^{\mathrm{firm}}-h^{\mathrm{soc}},
\]
using $h^{\mathrm{firm}}>h^{\mathrm{soc}}$. Removing absolute values:
\[
-(h^{\mathrm{firm}}-h^{\mathrm{soc}})<h^{\mathrm{soc}}-\bar h<h^{\mathrm{firm}}-h^{\mathrm{soc}}
\;\Longleftrightarrow\;
2h^{\mathrm{soc}}-h^{\mathrm{firm}}<\bar h<h^{\mathrm{firm}}.
\]
Maximising \eqref{eq:dW} in $\bar h$ gives the peak $\bar h=h^{\mathrm{soc}}$ with gain
\[
\Delta W(h^{\mathrm{soc}})=\tfrac{\alpha+\beta}{2}\bigl(h^{\mathrm{firm}}-h^{\mathrm{soc}}\bigr)^{2}=\tfrac{4}{2}(0.10)^{2}=0.020.
\]
\textbf{Band collapse at $\alpha=\beta$.} Substituting the bliss-point formulas,
\[
2h^{\mathrm{soc}}-h^{\mathrm{firm}}=\frac{2(A-\tau w)}{\alpha+\beta}-\frac{A-w}{\beta}\xrightarrow{\alpha=\beta}\frac{A-\tau w}{\beta}-\frac{A-w}{\beta}=\frac{(1-\tau)w}{\beta}=h^{*}.
\]
At $\alpha=\beta$ the lower band-edge equals the worker's bliss $h^{*}$, so the band simplifies to $\bar h\in(h^{*},h^{\mathrm{firm}})=(0.30,0.50)$.

\subsection{Worker-surplus band and the two thresholds}
The sign of $\Delta V_w$ comes from \eqref{eq:dVw}. For $\bar h<h^{\mathrm{firm}}$, the leading factor $(\bar h-h^{\mathrm{firm}})<0$, so $\Delta V_w(\bar h)>0$ iff the bracket is also negative:
\[
(1-\tau)w<\tfrac{\alpha}{2}(\bar h+h^{\mathrm{firm}})
\;\Longleftrightarrow\;
\bar h+h^{\mathrm{firm}}>\frac{2(1-\tau)w}{\alpha}=2h^{*}
\;\Longleftrightarrow\;
\bar h>2h^{*}-h^{\mathrm{firm}}.
\]
Combined with $\bar h<h^{\mathrm{firm}}$ (cap binding), the worker-surplus band is
\[
\Delta V_w(\bar h)>0\;\Longleftrightarrow\;\bar h\in\bigl(2h^{*}-h^{\mathrm{firm}},\,h^{\mathrm{firm}}\bigr).
\]
The peak of $\Delta V_w$ is at the worker's bliss $\bar h=h^{*}$, with magnitude
\[
\Delta V_w(h^{*})=V_w(h^{*})-V_w(h^{\mathrm{firm}})=\tfrac{\alpha}{2}(h^{\mathrm{firm}}-h^{*})^{2}=\tfrac{2}{2}(0.20)^{2}=0.04
\]
at calibration. Outside the band ($\bar h<2h^{*}-h^{\mathrm{firm}}$), the cap is so tight that the worker is forced past her bliss far enough to be worse off than at $h^{\mathrm{firm}}$.

\textbf{The two thresholds reconciled.} The joint-surplus and worker-surplus bands have different lower edges:
\[
\underbrace{2h^{\mathrm{soc}}-h^{\mathrm{firm}}}_{\text{joint turns negative}}
\;\ge\;
\underbrace{2h^{*}-h^{\mathrm{firm}}}_{\text{worker turns negative}}
\qquad(\text{since }h^{\mathrm{soc}}\ge h^{*}).
\]
Equality holds at the no-tax limit $\tau\!=\!0$ (where $h^{\mathrm{soc}}\!=\!h^{*}$ in this calibration); generically the joint threshold is strictly above the worker threshold. At baseline $(2h^{\mathrm{soc}}-h^{\mathrm{firm}},\,2h^{*}-h^{\mathrm{firm}})=(0.30,\,0.10)$, so for $\bar h\in(0.10,0.30)$ the cap is welfare-reducing on aggregate (firm losses dominate) but the worker individually still benefits.

% =========================================================================
\section{Formal derivation of the Q5 RTM bargain}
% =========================================================================

\subsection{Composite payoffs along the labour-demand curve}
Per sector $j\in\{\textsc{g},\textsc{b}\}$, the firm responds on its labour-demand curve $h^d_j(w)=(A_j-w)/\beta$. Substituting into the primitives:
\begin{align*}
V_w(w;A_j) &\equiv V_w(w,h^d_j(w))=\frac{(1-\tau)w(A_j-w)}{\beta}-\frac{\alpha(A_j-w)^{2}}{2\beta^{2}},\\
V_f(w;A_j) &\equiv V_f(w,h^d_j(w);A_j)=\frac{(A_j-w)^{2}}{2\beta},
\end{align*}
where $V_f$ collapses to the standard Cournot-residual form $(A_j-w)^{2}/(2\beta)$ because the firm is on its own FOC. Reparameterise $x_j\equiv A_j-w=\beta h^d_j(w)$; then
\begin{equation}
V_w(x_j)=A_v\,x_j-B_v\,x_j^{2},\qquad A_v\equiv\frac{(1-\tau)A_j}{\beta},\quad B_v\equiv\frac{1-\tau}{\beta}+\frac{\alpha}{2\beta^{2}};\qquad V_f(x_j)=\frac{x_j^{2}}{2\beta}.
\label{eq:rtm-composite}
\end{equation}

\subsection{Lagrangian for the generalised Nash programme}
Union and firm choose $w$ to maximise the asymmetric Nash product subject to the participation (IR) constraints:
\[
\max_{w}\;\bigl[V_w(w;A_j)-V_w^{0}\bigr]^{\eta}\,V_f(w;A_j)^{1-\eta}\qquad\text{s.t.}\quad V_w\ge V_w^{0},\ V_f\ge 0.
\]
Taking logs (both factors strictly positive on the interior), the Lagrangian with multipliers $\mu_w,\mu_f\ge 0$ is
\[
\mathcal{L}(x_j;\mu_w,\mu_f)
=\eta\ln\!\bigl(V_w(x_j)-V_w^{0}\bigr)+(1-\eta)\ln V_f(x_j)+\mu_w\bigl(V_w(x_j)-V_w^{0}\bigr)+\mu_f V_f(x_j),
\]
using $w=A_j-x_j$ and $dw/dx_j=-1$ (so the FOC in $x_j$ is the negative of the FOC in $w$, leaving the zero-locus unchanged). On the interior, $\mu_w=\mu_f=0$ and the FOC is
\begin{equation}
\frac{\eta\,V_w'(x_j)}{V_w(x_j)-V_w^{0}}+\frac{(1-\eta)\,V_f'(x_j)}{V_f(x_j)}=0,\qquad V_w'=A_v-2B_vx_j,\ V_f'=x_j/\beta.
\label{eq:rtm-foc}
\end{equation}

\subsection{Reduction to a quadratic in $x_j$}
Multiply \eqref{eq:rtm-foc} through by $\bigl(V_w-V_w^{0}\bigr)V_f$, substitute $V_f=x_j^{2}/(2\beta)$, and cancel one factor of $x_j/(2\beta)$ (assuming $x_j\!\neq\!0$, which excludes the no-employment corner $w=A_j$):
\[
\eta\,(A_v-2B_vx_j)\,x_j+2(1-\eta)\,\bigl(A_vx_j-B_vx_j^{2}-V_w^{0}\bigr)=0.
\]
Collecting terms:
\begin{equation}
2B_v\,x_j^{2}-A_v(2-\eta)\,x_j+2(1-\eta)\,V_w^{0}=0.
\label{eq:rtm-quad}
\end{equation}
The equilibrium $x_j^{\star}(\eta)$ is the larger root (the smaller root is the IR-binding fixed point at the $\eta\!=\!0$ corner --- see below), and the bargained wage and hours follow as $w_j(\eta)=A_j-x_j^{\star}(\eta)$, $h_j(\eta)=x_j^{\star}(\eta)/\beta$. At calibration $(\alpha,\beta,\tau,V_w^{0})=(2,2,0.4,0.05)$, \eqref{eq:rtm-quad} reads $1.1\,x^{2}-0.3A_j(2-\eta)x+0.1(1-\eta)=0$. Two checks: at $(A,\eta)=(2,0)$ we get $x^{\star}=1$, recovering $w=1,\,h=0.5$; at $(A_g,\eta)=(2.2,0.5)$ we get $x^{\star}\!\approx\!0.846$, giving $w\!\approx\!1.354$, $h\!\approx\!0.423$.

\subsection{Boundary conditions}
\textbf{(i) $\eta=0$.} Equation \eqref{eq:rtm-quad} becomes $2B_vx_j^{2}-2A_vx_j+2V_w^{0}=0$, i.e.\ $V_w(x_j)=V_w^{0}$. The worker's IR binds and the firm picks the wage that strips all rent above the outside option; on the firm's labour-demand curve this delivers $h_j=h^{\mathrm{firm}}_j$.

\textbf{(ii) $\eta\to 1$.} \eqref{eq:rtm-quad} becomes $2B_vx_j^{2}-A_vx_j=0$, with non-trivial root
\[
x_j^{\star}(1)=\frac{A_v}{2B_v}=\frac{(1-\tau)A_j\,\beta}{2\beta(1-\tau)+\alpha}\;\Longrightarrow\;
h_j(1)=\frac{x_j^{\star}(1)}{\beta}=\frac{(1-\tau)A_j}{2\beta(1-\tau)+\alpha}.
\]
\emph{SOC at the limit.} Evaluating the bracket $4B_v x_j-A_v(2-\eta)$ at $(x,\eta)=(A_v/(2B_v),1)$ gives $4B_v\!\cdot\!A_v/(2B_v)-A_v\!=\!2A_v-A_v\!=\!A_v>0$, so the second-order condition for the Nash maximum holds at the limit and $x_j^{\star}(1)$ is indeed the maximiser, not a corner.

\subsection{Walrasian benchmark and the comparative statics}
\textbf{Maintained regularity.} Two conditions on $V_w^{0}$ are imposed throughout. \emph{(R1)} $A_v^{2}\!>\!4B_vV_w^{0}$ (strict discriminant), so the IR-binding equation $V_w(x_j)\!=\!V_w^{0}$ admits two distinct real roots. \emph{(R2)} $x_j^{\mathrm{FM}}(A_j)^{2}>2\beta^{2}V_w^{0}/\alpha$ over the range of $A_j$ considered, which delivers $\partial w_j^{\mathrm{FM}}/\partial A_j>0$. At calibration: (R1) $A_v^{2}\!=\!0.36>4B_vV_w^{0}\!=\!4(0.55)(0.05)\!=\!0.11$. (R2) $x_j^{\mathrm{FM}}\!=\!1$ and $2\beta^{2}V_w^{0}/\alpha\!=\!0.2$, so (R2) holds with slack; on the sectoral range $A_j\!\in\![1.8,2.2]$ used in Q5 it remains satisfied (smallest value $x_b^{\mathrm{FM}}\!=\!0.878$ gives $x^{2}\!=\!0.771>0.2$).

The competitive (free-market) outcome is the $\eta\!=\!0$ specialisation of \eqref{eq:rtm-quad}: $V_w(x_j)\!=\!V_w^{0}$, with larger root
\begin{equation}
x_j^{\mathrm{FM}}(A_j)=\frac{A_v+\sqrt{A_v^{2}-4B_v V_w^{0}}}{2B_v},\qquad h_j^{\mathrm{FM}}=\frac{x_j^{\mathrm{FM}}}{\beta},\qquad w_j^{\mathrm{FM}}=A_j-x_j^{\mathrm{FM}}.\label{eq:walras}
\end{equation}
Differentiating $V_w(x_j;A_j)\!=\!V_w^{0}$ implicitly and using $\partial A_v/\partial A_j=(1-\tau)/\beta$:
\[
\frac{\partial x_j^{\mathrm{FM}}}{\partial A_j}=\frac{(1-\tau)\,x_j^{\mathrm{FM}}}{\beta(2B_vx_j^{\mathrm{FM}}-A_v)}>0,
\]
positive because the larger root sits past the vertex $x\!=\!A_v/(2B_v)$ of the parabola $V_w$, where $V_w'\!=\!A_v\!-\!2B_vx<0$. So $\partial h_j^{\mathrm{FM}}/\partial A_j\!=\!\beta^{-1}\partial x_j^{\mathrm{FM}}/\partial A_j>0$. For the wage, substituting $A_v x_j\!=\!B_v x_j^{2}\!+\!V_w^{0}$ into the numerator above gives $\partial x_j^{\mathrm{FM}}/\partial A_j\!=\!(1-\tau)x_j^{2}/[\beta(B_vx_j^{2}-V_w^{0})]<1$ whenever $x_j^{2}\!>\!2\beta^{2}V_w^{0}/\alpha$ (using $\beta B_v\!-\!(1-\tau)\!=\!\alpha/(2\beta)$); this is exactly condition (R2), so $\partial w_j^{\mathrm{FM}}/\partial A_j\!=\!1-\partial x_j^{\mathrm{FM}}/\partial A_j>0$. Together,
\[
\frac{\partial h_j^{\mathrm{FM}}}{\partial A_j}>0,\qquad \frac{\partial w_j^{\mathrm{FM}}}{\partial A_j}>0,
\]
i.e.\ ``hours follow productivity'' under perfect labour mobility.

\subsection{Comparative statics}
Differentiating \eqref{eq:rtm-quad} implicitly in $\eta$:
\[
\bigl[4B_vx_j-A_v(2-\eta)\bigr]\frac{dx_j}{d\eta}=-A_vx_j+2V_w^{0}.
\]
\textbf{Bracket positivity (global on $[0,1)$).} The larger root of \eqref{eq:rtm-quad} is $x_j^{\star}(\eta)\!=\![A_v(2-\eta)+\sqrt{D(\eta)}]/(4B_v)$ with $D(\eta)\!\equiv\!A_v^{2}(2-\eta)^{2}-16B_v(1-\eta)V_w^{0}$. At $\eta\!=\!0$, $D(0)\!=\!4A_v^{2}-16B_vV_w^{0}\!=\!4(A_v^{2}-4B_vV_w^{0})>0$ by (R1); at $\eta\!\to\!1$, $D(1)\!=\!A_v^{2}\!>\!0$. $D(\eta)$ is a quadratic in $\eta$ with positive leading coefficient $A_v^{2}$, vertex at $\eta_v\!=\!2(1-4B_vV_w^{0}/A_v^{2})$, and minimum value $D_{\min}\!=\!16B_vV_w^{0}(A_v^{2}-4B_vV_w^{0})/A_v^{2}>0$ under (R1). At calibration $\eta_v\!=\!1.39$ lies outside $[0,1]$, so $D$ is monotone on the relevant interval (decreasing from $D(0)\!=\!1.0$ to $D(1)\!=\!0.36$). Therefore $4B_vx_j^{\star}(\eta)\!=\!A_v(2-\eta)+\sqrt{D(\eta)}>A_v(2-\eta)$ on $[0,1)$, so the bracket $4B_vx_j-A_v(2-\eta)\!=\!\sqrt{D(\eta)}>0$ strictly. The bracket equals $-V_w'(x_j^{\star})\cdot(\text{positive})$ at the larger root, where $V_w'(x_j^{\star})\!<\!0$ (root past the vertex), so positivity of the bracket is the second-order condition for the Nash maximum.

\textbf{Sign of the right-hand side.} $-A_vx_j+2V_w^{0}<0$ iff $A_v x_j>2V_w^{0}$. At the larger root we have $x_j^{\star}(\eta)\!\ge\!x_j^{\star}(1)\!=\!A_v/(2B_v)$, so $A_v x_j^{\star}\!\ge\!A_v^{2}/(2B_v)>2V_w^{0}$ where the last inequality is exactly (R1). So the RHS is strictly negative on $[0,1]$, and combined with the strictly positive bracket:
\[
\frac{dx_j}{d\eta}<0\;\Longrightarrow\;\frac{dw_j}{d\eta}=-\frac{dx_j}{d\eta}>0,\qquad \frac{dh_j}{d\eta}=\frac{1}{\beta}\frac{dx_j}{d\eta}<0,
\]
the union-pulls-wages-up / hours-down comparative statics.

\subsection{The welfare-restorer claim}
\textbf{Joint surplus.} $W$ is concave in $h$ with peak at $h^{\mathrm{soc}}_j\in(h^{*}_j,h^{\mathrm{firm}}_j)$. Since $h_j(\eta)$ is continuous and strictly decreasing in $\eta$ with $h_j(0)=h^{\mathrm{firm}}_j$, by the intermediate-value theorem there exists a unique $\hat\eta_j\in(0,1)$ with $h_j(\hat\eta_j)=h^{\mathrm{soc}}_j$. Combining with the quadratic identity \eqref{eq:dW},
\[
W\bigl(h_j(\eta)\bigr)-W(h^{\mathrm{firm}}_j)=-\tfrac{\alpha+\beta}{2}\Bigl[(h^{\mathrm{soc}}_j-h_j(\eta))^{2}-(h^{\mathrm{soc}}_j-h^{\mathrm{firm}}_j)^{2}\Bigr],
\]
which is increasing on $[0,\hat\eta_j]$ (the union pulls hours \emph{toward} the social bliss) and decreasing on $[\hat\eta_j,1]$ (overshoot region).

\textbf{Worker surplus is monotone in $\eta$.} On the labour-demand curve $V_w(x_j)=A_v x_j-B_vx_j^{2}$ is a concave parabola with maximum at $\hat x\!\equiv\!A_v/(2B_v)\!=\!x_j^{\star}(1)$. At $\eta\!=\!0$, $V_w(x_j^{\star}(0))=V_w^{0}$, so $x_j^{\star}(0)$ is the larger root of $V_w(x)=V_w^{0}$:
\[
x_j^{\star}(0)=\frac{A_v+\sqrt{A_v^{2}-4B_vV_w^{0}}}{2B_v}>\frac{A_v}{2B_v}=\hat x=x_j^{\star}(1),
\]
which together with monotonicity of $x_j^{\star}(\eta)$ gives $x_j^{\star}(\eta)\in(\hat x,x_j^{\star}(0)]$ for all $\eta\in[0,1)$. On this interval $V_w'(x)=A_v-2B_vx<0$, so as $x_j^{\star}$ falls in $\eta$, $V_w$ rises:
\[
\frac{dV_w}{d\eta}=V_w'(x_j^{\star})\cdot\frac{dx_j^{\star}}{d\eta}>0\qquad\text{for all }\eta\in[0,1),
\]
a product of two negatives.

\textbf{Conclusion.} Joint surplus is hump-shaped in $\eta$ on $[0,1]$ with peak at $\hat\eta_j$, and worker surplus is strictly increasing in $\eta$. A planner with any positive worker weight $\lambda>0$ in the Q7 weighting $W^{\lambda}(\eta)=(1+\lambda)V_w+(1-\lambda)V_f$ thus strictly prefers $\eta\!>\!0$ to $\eta\!=\!0$, since $V_f$ is bounded above by its $\eta\!=\!0$ value while $V_w$ is strictly higher.

\end{document}
